% ===================================================================
%  TABEL Nr. 5 — Het huis en de bouw ervan.
%  Traduction 2026 d'apres texto/io, controlee sur texto/fr et
%  texto/es. Consigne : voir texto/nl/10-tabel-01.tex.
%
%  Le tableau 5 est un DIALOGUE, et les deux volumes composent le nom
%  du parleur en petites capitales : on garde \textsc{}.
%
%  LA PREMIERE INVERSION DE LA COLONNE, et elle est de syntaxe. L'ido
%  ecrit « il indutas per gipso (70) la plafoni (26) » — l'instrument
%  avant l'objet ; le neerlandais ordinaire renverse : « hij
%  bepleistert de plafonds met gips », et le releve sortait 26 avant
%  70. On ne deplace pas le renvoi, qui appartient a son mot : on met
%  le complement en tete — « Met gips (70) bepleistert hij de plafonds
%  (26) en de muren » — ce que le neerlandais fait tous les jours,
%  parce que son verbe reste en deuxieme position quoi qu'on place
%  devant lui. L'ordre est celui de l'ido, et la phrase reste
%  neerlandaise.
%
%  LES DEUX NOMS PROPRES SE TRADUISENT, ET L'IDO L'A DEJA FAIT :
%  Leblanc est « Blanko », Roux est « Rufo ». Le neerlandais forme les
%  memes noms sur les memes adjectifs — wit, ros — et ecrit donc
%  mijnheer De Wit et mijnheer De Rooij, qui sont de vrais noms de
%  famille neerlandais.
% ===================================================================

%%K t05-apar-1 apar
\VUsaut{6.0mm}
\VUtitre{15.0pt}{60}{TABEL Nr. 5}
\VUsaut{1.4mm}
\VUtitre{11.4pt}{0}{Het huis en de bouw ervan.}
\VUsaut{2.2mm}

%%K t05-tit-1 sub
\VUpk{10.2pt}{40}{Een gelukkige ontmoeting.}
\VUsaut{3.0mm}

%%K t05-01-1 p
\textsc{Jan}. --- Oom, ik zou heel graag willen weten hoe men een \VUgras{huis} bouwt.

%%K t05-01-2 p
\textsc{De oom} (80). --- Dat is heel gemakkelijk, mijn jongen, kom met mij mee, ik zal je het
huis laten zien dat mijnheer Bernard (79) laat bouwen en dat ik zal bewonen.

%%K t05-01-3 p
\textsc{Jan}. --- Wie heeft dan het \VUgras{plan} \textsuperscript{(76)} van dit huis
getekend?

%%K t05-01-4 p
\textsc{De oom}. --- \VUgras{De architect} \textsuperscript{(75)}; hij heeft een heel
duidelijke tekening voorgelegd die goedgekeurd werd, en \VUgras{de aannemer}
\textsuperscript{(77)} is met de werken begonnen.

%%K t05-01-5 p
\textsc{Jan}. --- Wie is dan de aannemer?

%%K t05-01-6 p
\textsc{De oom}. --- Hij is mijnheer De Wit, de vader van je vriend Jacob. Hij leidt de
werklieden samen met mijnheer De Rooij, de architect.

%%K t05-01-7 p
Kijk! juist op tijd komt mijnheer De Wit met zijn \VUgras{duimstok} \textsuperscript{(78)}, en
mijnheer De Rooij met zijn plan.

%%K t05-01-8 p
Goedendag, heren. Hoe gaat het met u?

%%K t05-01-9 p
\textsc{Mijnheer De Rooij}. --- Heel goed, dank u. Goedendag, Jan, wat is dat kind gegroeid!

%%K t05-01-10 p
\textsc{De oom}. --- Ja, werkelijk, hij is een kleine man. Hij is heel ernstig, men moet hem
over alles wat hij ziet inlichten. Vandaag zou hij willen weten hoe men een huis bouwt.

%%K t05-tit-2 sub
\VUpk{10.2pt}{40}{De werklieden.}
\VUsaut{3.0mm}

%%K t05-01-11 p
\textsc{Mijnheer De Wit}. --- Wel, kom met mij mee, vriendje, ik zal het je dadelijk
uitleggen, want ik houd veel van kinderen die zich willen onderrichten.

%%K t05-01-12 p
Je ziet die werkman die een \VUgras{schop} \textsuperscript{(82)} in de hand houdt: hij is een
\VUgras{grondwerker} \textsuperscript{(81)}. Hij graaft een
\VUgras{sleuf} \textsuperscript{(46)} om de waterleiding te leggen. Hij heeft
ook de grond uitgegraven op de plaats waar het huis staat, opdat de metselaar de vier
\VUgras{muren} zou optrekken. Het deel van die muren dat in de grond zit wordt
\VUgras{fundering} \textsuperscript{(2)} genoemd. De ruimte die tussen de funderingen ligt
vormt de \VUgras{kelder} \textsuperscript{(3)}. Die kelder heeft een klein raam dat
\VUgras{keldergat} \textsuperscript{(5)} heet, een
\VUgras{deur} \textsuperscript{(4)} en een trap.

%%K t05-01-13 p
De \VUgras{metselaar} \textsuperscript{(108)} is doorgegaan met het bouwen van de muren en
heeft openingen gelaten voor de \VUgras{deuren} \textsuperscript{(8)} en de \VUgras{ramen}
\textsuperscript{(17)}.

%%K t05-01-14 p
\textsc{Jan}. --- Maar die metselaar zie ik niet!

%%K t05-01-15 p
Mijnheer \textsc{De Wit}. --- Nu is hij klaar met werken in het huis. Hij staat bij de
\VUgras{paardenstal} \textsuperscript{(54)}, op zijn \VUgras{steiger} \textsuperscript{(110)},
hij bouwt de muur van het \VUgras{washuis} \textsuperscript{(56)}.

%%K t05-01-16 p
Hij heeft een truweel, een kuip en een \VUgras{schietlood} \textsuperscript{(109)}. Maar die
namen zul je niet onthouden.

%%K t05-01-17 p
\textsc{Jan}. --- Zeker wel, want ik zal dadelijk aan mijn oom een klein truweel en een kuip
vragen, en ik zal zelf een schietlood met een touwtje maken.

%%K t05-tit-3 sub
\VUsaut{3.0mm}
\VUpk{10.2pt}{40}{Het materiaal.}
\VUsaut{3.0mm}

%%K t05-01-18 p
\textsc{De oom}. --- Ha! heel goed. Maar zeg mij eens, Jan, wat heeft men nodig om een huis te
bouwen?

%%K t05-01-19 p
\textsc{Jan}. --- Men heeft \VUgras{hardsteen} \textsuperscript{(59)} en
\VUgras{mortel} \textsuperscript{(65)} nodig.

%%K t05-01-20 p
\textsc{De oom}. --- Juist, zie die man, met een grote hoed, op zijn
\VUgras{kar} \textsuperscript{(136)}. Hij is een \VUgras{voerman}
\textsuperscript{(135)}. Elke dag brengt hij hier \VUgras{steenblokken}
\textsuperscript{(60)}. Hij laadt zijn wagen af op de binnenplaats, en de
\VUgras{steenhouwers} \textsuperscript{(83)} bekappen de blokken en zagen ze met
hun \VUgras{steenzaag} \textsuperscript{(85)}. Een \VUgras{opperman} \textsuperscript{(146)}
draagt ze naar de metselaar die ze op hun plaats legt.

%%K t05-01-21 p
Intussen maakt de \VUgras{mortelmaker} \textsuperscript{(87)} de mortel klaar. Hij heeft
\VUgras{zand} \textsuperscript{(62)}, \VUgras{kalk} \textsuperscript{(63)} en \VUgras{water}
\textsuperscript{(64)} nodig. De kalk staat naast hem in een \VUgras{zak}
\textsuperscript{(71)}, een
\VUgras{metselaarsknecht} \textsuperscript{(111)} brengt hem het zand op een
\VUgras{kruiwagen} \textsuperscript{(112)} en de \VUgras{leerjongen}
\textsuperscript{(113)} staat bij de pomp (58). Hij vult een \VUgras{emmer}
\textsuperscript{(114)}. De mortel is weldra klaar. De \VUgras{morteldrager}
\textsuperscript{(107)} neemt hem en draagt hem langs de ladder omhoog, op de steiger, waar
een metselaar werkt die die zijde van het huis afmaakt.

%%K t05-01-22 p
En nu, hoe wordt de werkman genoemd die het \VUgras{dak} \textsuperscript{(31)} maakt?

%%K t05-01-23 p
\textsc{Jan}. --- Hij is een \VUgras{dakdekker} \textsuperscript{(118)}.

%%K t05-tit-4 sub
\VUsaut{3.0mm}
\VUpk{10.2pt}{40}{Het dak van het huis.}
\VUsaut{3.0mm}

%%K t05-01-24 p
Mijnheer \textsc{De Wit}. --- Ja, maar eerst komt de \VUgras{timmerman}
\textsuperscript{(115)}.

%%K t05-01-25 p
De timmerman maakt het \VUgras{dakgebint} \textsuperscript{(35)}. Hij verbindt de
\VUgras{balken} \textsuperscript{(37)} met \VUgras{houten pennen} \textsuperscript{(117)}. Die
werklieden, daarboven, met hun wijde fluwelen broeken, zijn timmerlieden. De een houdt een
balkje vast, en de ander zaagt een pen af met zijn \VUgras{bijl} \textsuperscript{(116)}.
Degene die bij de
\VUgras{schoorsteen} \textsuperscript{(41)} staat is de dakdekker. Hij spijkert
latten op de (*) balken, en \VUgras{leien} \textsuperscript{(66)} op de latten. Soms legt hij
dakpannen.

%%K t05-noto-1 noto
\VUnotes{143.2mm}{%
(*) \VUgras{Balk-o} = D. \textit{Balken}, E. \textit{joist}, F. \textit{solive}, I.
\textit{travicello}, R. \textit{balka}, S. Port. \textit{viga}. Technische wortel die tot nu
toe niet aangenomen is, maar voor te stellen om de zin van het F. woord \textit{solive} te
vertalen, dat noch een balk F. \textit{poutre}, noch een balkje is. Het is een vierkant
gemaakte houten stang die een vloer en een plafond draagt.%
}

%%K t05-01-26 p
Het dak van de paardenstal is met \VUgras{dakpannen gedekt} \textsuperscript{(55)}, dat van
het huis met \VUgras{leien} \textsuperscript{(32)} en dat van de veranda is van \VUgras{zink}
\textsuperscript{(24)}.

%%K t05-01-27 p
\textsc{Jan}. --- Maakt de dakdekker ook de zinken daken?

%%K t05-01-28 p
Mijnheer \textsc{De Wit}. --- Nee, maar de \VUgras{zinkwerker} \textsuperscript{(121)} of de
\VUgras{loodgieter}, degene die een
\VUgras{dakvenster} \textsuperscript{(38)} op het dak van het huis plaatst. Hij
plaatst ook de \VUgras{regenpijpen} \textsuperscript{(43)} en de
\VUgras{dakgoten} \textsuperscript{(42)}. Hij soldeert het zink en het blik. Hij
heeft de \VUgras{bliksemafleider} \textsuperscript{(40)} en de
\VUgras{windwijzer} \textsuperscript{(39)} op de \VUgras{geveltop}
\textsuperscript{(36)} bevestigd.

%%K t05-01-29 p
\textsc{Jan}. --- Is het huis dan afgebouwd?

%%K t05-tit-5 sub
\VUsaut{3.0mm}
\VUpk{10.2pt}{40}{Het gereedschap.}
\VUsaut{3.0mm}

%%K t05-01-30 p
Mijnheer \textsc{De Wit}. --- Niet helemaal. De \VUgras{stukadoor} \textsuperscript{(99)}
bouwt met \VUgras{bakstenen} \textsuperscript{(61)} de scheidingswanden die het in
verschillende kamers verdelen. Met \VUgras{gips} \textsuperscript{(70)} bepleistert hij de
\VUgras{plafonds} \textsuperscript{(26)} en de muren. Nu bewerkt hij het plafond van de
\VUgras{veranda} \textsuperscript{(33)}. Hij heeft, zoals de metselaar, een \VUgras{truweel}
\textsuperscript{(100)} en een \VUgras{kuip} \textsuperscript{(101)}.

%%K t05-01-31 p
Daarna maakt de schrijnwerker de deuren, de ramen, de \VUgras{parketvloeren}
\textsuperscript{(16)} en de \VUgras{blinden} \textsuperscript{(19)}.

%%K t05-01-32 p
Nu werkt hij op de binnenplaats aan zijn \VUgras{werkbank} \textsuperscript{(90)}. Hij heeft
een \VUgras{zaag} \textsuperscript{(94)} om het hout (69) te zagen, een \VUgras{schaaf}
\textsuperscript{(91)}, een
\VUgras{klemschroef} \textsuperscript{(92)}, een \VUgras{beitel}
\textsuperscript{(94 bis)}, een \VUgras{passer} \textsuperscript{(95)}, en een houten
\VUgras{hamer} \textsuperscript{(95 bis)}.

%%K t05-01-33 p
\textsc{Jan}. --- Ha! maar schrijnwerken is leuker dan metselen. Oom, u zult voor mij een
werkbank, een zaag en een schaaf kopen, want ik zal weldra een hok voor Medor maken.

%%K t05-01-34 p
\textsc{De oom}. --- Ja, jongen.

%%K t05-01-35 p
\textsc{Jan}. --- Wat ben ik blij! En wie is degene die bij de voordeur staat?

%%K t05-tit-6 sub
\VUsaut{3.0mm}
\VUpk{10.2pt}{40}{Het schilderen.}
\VUsaut{3.0mm}

%%K t05-01-36 p
Mijnheer \textsc{De Wit}. --- Hij is een \VUgras{schilder} \textsuperscript{(102)}. Hij staat
op zijn ladder met een pot \VUgras{verf} \textsuperscript{(103)} en zijn \VUgras{kwast}
\textsuperscript{(104)}. Hij schildert het hout van de voordeur om het langer te bewaren.

%%K t05-01-37 p
Hij plakt ook het behang in de vertrekken.

%%K t05-01-38 p
De \VUgras{stoffeerder} \textsuperscript{(107)} die op een \VUgras{ladder}
\textsuperscript{(106)} staat, op het \VUgras{balkon} \textsuperscript{(28)}, brengt een
\VUgras{zonnescherm} \textsuperscript{(29)} aan.

%%K t05-01-39 p
De werkman die bij het grote raam staat is de \VUgras{glazenmaker} \textsuperscript{(96)}. Hij
zet de \VUgras{ruiten} \textsuperscript{(18)} vast met \VUgras{stopverf}
\textsuperscript{(73)}. Die andere werkman, geknield bij de kelderdeur is een
\VUgras{slotenmaker} \textsuperscript{(97)}. Hij plaatst een \VUgras{slot}
\textsuperscript{(6)}. Zijn \VUgras{vijl} \textsuperscript{(98)} ligt naast hem.

%%K t05-01-40 p
\textsc{Jan}. --- Is die werkman die met zijn \VUgras{hamer} \textsuperscript{(84)} op een
ijzerwerk slaat ook een slotenmaker?

%%K t05-01-41 p
Mijnheer \textsc{De Wit}. --- Juister gezegd, hij is een smid (125). Hij smeedt
\VUgras{ijzerwerk} \textsuperscript{(67)} voor de \VUgras{elektricien}
\textsuperscript{(124)} die op het dak van het \VUgras{koetshuis} \textsuperscript{(51)}
staat. Hij heeft een \VUgras{smidse} \textsuperscript{(126)}, een \VUgras{aambeeld}
\textsuperscript{(127)} en een
\VUgras{tang} \textsuperscript{(128)}.

%%K t05-01-42 p
De elektricien bevestigt de \VUgras{koperen} \textsuperscript{(44)}
\VUgras{draad} van de telefoon.

%%K t05-tit-7 sub
\VUpk{10.2pt}{40}{Het tuinieren.}
\VUsaut{3.0mm}

%%K t05-01-43 p
\textsc{De oom}. --- Achterin de \VUgras{binnenplaats} \textsuperscript{(45)}, bij het
\VUgras{hek} \textsuperscript{(47)} en de
\VUgras{poort} \textsuperscript{(48)} zie je de \VUgras{tuinman}
\textsuperscript{(129)}. Hij maakt het \VUgras{tuintje} \textsuperscript{(49)} klaar. Hij
heeft de grond omgespit met zijn \VUgras{spade} \textsuperscript{(131)}, hij zaait zaden en
harkt met de \VUgras{hark} \textsuperscript{(130)}. Daarna zal hij met zijn \VUgras{gieter}
\textsuperscript{(132)} de \VUgras{bloembedden} \textsuperscript{(150)} begieten en de planten
zullen groeien.

%%K t05-01-44 p
De \VUgras{stalknecht} \textsuperscript{(133)} die zojuist het paard verzorgd en het voer
gegeven heeft maakt het \VUgras{rijtuig} \textsuperscript{(52)} schoon met een spons, een
\VUgras{handpomp} \textsuperscript{(134)} en een emmer water. Daarna zal hij het weer in het
koetshuis rijden en de deur sluiten met de dikke
\VUgras{grendel} \textsuperscript{(53)}.

%%K t05-01-45 p
Ziezo, je bent tevreden, denk ik, je hebt genoeg uitleg gekregen.

%%K t05-01-46 p
\textsc{Jan}. --- Ja, oom, maar ik zou graag de binnenkant van het huis zien.

%%K t05-01-47 p
\textsc{De oom}. --- Dat zal morgen gebeuren. Mijnheer De Rooij zal je het plan uitleggen en
je de naam van elke kamer noemen.

%%K t05-tit-8 sub
\VUsaut{3.0mm}
\VUpk{10.2pt}{40}{Inwendige indeling van het huis.}
\VUsaut{2.4mm}

%%K t05-apar-2 apar
{\centering\textit{(Zie het plan.)}\par}
\VUsaut{2.4mm}

%%K t05-01-48 p
\textsc{De architect}. --- Kom, Jan, ik zal je dadelijk de verschillende delen van het huis
laten verkennen.

%%K t05-01-49 p
Laten wij de \VUgras{begane grond} \textsuperscript{(7)} binnengaan. Dit is de knop van de
\VUgras{schel} \textsuperscript{(11)}. Wij gaan de drie
\VUgras{treden} \textsuperscript{(14)} van het \VUgras{bordes}
\textsuperscript{(9)} op, wij overschrijden de \VUgras{drempel} \textsuperscript{(10)} van de
\VUgras{voordeur} \textsuperscript{(8)} en dan staan wij aan de voet van de \VUgras{trap}
\textsuperscript{(13)}.

%%K t05-01-50 p
\textsc{Jan}. --- Dat is de gang.

%%K t05-01-51 p
\textsc{De architect}. --- Nee, dat is de \VUgras{vestibule} \textsuperscript{(12)}. De
\VUgras{gang} \textsuperscript{(a)} is aan het uiteinde. Rechts liggen de \VUgras{salon}
\textsuperscript{(b)} en de
\VUgras{eetkamer} \textsuperscript{(c)}, tegenover de eetkamer zijn de
\VUgras{keuken} \textsuperscript{(d)} en de \VUgras{dienstkamer}
\textsuperscript{(e)}, daarna vooraan het \VUgras{studeervertrek} \textsuperscript{(f)}. De
\VUgras{veranda} \textsuperscript{(23)} met haar
\VUgras{glazen deur} \textsuperscript{(25)} ligt naast de salon.

%%K t05-01-52 p
Wij zullen dadelijk de trap opgaan. Houd je vast aan de \VUgras{leuning}
\textsuperscript{(15)} en tel de treden. Het zijn er veertien. We staan op de
\VUgras{overloop} \textsuperscript{(m)}, wij zijn op de eerste
\VUgras{verdieping} \textsuperscript{(27)}.

%%K t05-tit-9 sub
\VUsaut{3.0mm}
\VUpk{10.2pt}{40}{De eerste verdieping.}
\VUsaut{3.0mm}

%%K t05-01-53 p
Tegenover de trap zijn het \VUgras{privaat} \textsuperscript{(20)} en het
\VUgras{toiletkamertje} \textsuperscript{(j)}. Rechts een mooie
\VUgras{slaapkamer} \textsuperscript{(g)} met twee ramen aan de \VUgras{gevel}
\textsuperscript{(1)} van het huis. Links zijn de \VUgras{badkamer} \textsuperscript{(1)} en
de \VUgras{logeerkamer} \textsuperscript{(i)} met een grote \VUgras{alkoof}
\textsuperscript{(h)}: Naast de slaapkamer is de
\VUgras{kinderkamer} \textsuperscript{(k)}. Zij grenst aan een ruim
\VUgras{terras} \textsuperscript{(21)}. Onder dat terras is nog een privaat (20)
en aan de muur hangt de \VUgras{bel} \textsuperscript{(22)} die voor het diner luidt.

%%K t05-01-54 p
\textsc{Jan}. --- Laten wij naar de \VUgras{tweede verdieping} gaan.

%%K t05-01-55 p
\textsc{De architect}. --- Er bestaat geen tweede verdieping. Onder het dak zijn alleen
\VUgras{mansardekamers} \textsuperscript{(33)} en de zolder (34). Wij zijn klaar met de
verkenning, wij kunnen weer naar beneden gaan.

%%K t05-01-56 p
\textsc{Jan}. --- Dank u, mijnheer, dat is heel boeiend. Ik zal dadelijk aan mijn vader alles
vertellen wat ik gezien heb.
\VUsaut{4.0mm}
\VUfilet{22mm}
